\begin{definition}
	Метод суммирования Чезаро. 
	\[\sum_{n=1}^\infty a_n,\ S_n=\sum_{k=1}^na_n,\ G_n:=\frac{S_1+\dots+S_n}{n}\]
	Если $\ex\lim_{n\ra\infty}G_n=S$, то $S$ называется суммой по Чезарю $\sum_{n=1}^\infty a_n$.
\end{definition}
\begin{theorem}\label{13_th1}
	Метод суммирования Чезаро регулярен, то есть $\fa$ сходящегося ряда $\sum_{n=1}^\infty a_n$ его сумма по Чезарю $\ex$ и равна обычной сумме.
\end{theorem}
\begin{proof}
	Пусть $L=\lim_{n\ra\infty}S_n$. Тогда:
	\[(\fa\eps>0)(\ex N\in\N)(\fa n>N)\ |S_n-L|<\frac\eps2\]
	\[G_n-L=\frac{S_1+\dots+S_n}n-\frac{nL}{n}=\frac{(S_1-L)+\dots+(S_N-L)}{n}+\frac{(S_{N+1}-L)+\dots+(S_n-L)}{n}=P_1+P_2\]
	Оценим $P_1$ и $P_2$:
	\[(\ex N_2\in\N)(\fa n>N_2)\ |P_1|<\frac\eps2\]
	\[|P_2|<\frac\eps2\cdot\frac{n-N_1}{n}<\frac\eps2\]
	Тогда получили:
	\[(\fa\eps>0)(\ex N=\max(N_1,N_2))(\fa n>N)\ |G_n-L|<\eps\lra\lim_{n\ra\infty}G_n=L\]
	Получили требуемое.
\end{proof}
\begin{note}
	Аналогично можно доказать, что если $\lim_{n\ra\infty}S_n=\pm\infty$, то
	$\lim_{n\ra\infty}G_n=\pm\infty$.
\end{note}
\begin{corollary}\label{13_cl1}
	Если $(\fa n\in\N)\ S_n>0,\ \lim_{n\ra\infty}S_n=L$, то $\lim_{n\ra\infty}\sqrt[n]{S_1\dots S_n}=L$.
\end{corollary}
\begin{proof}
	$\lim_{n\ra\infty}\ln S_n=\ln L$
\end{proof}
\begin{anote}
	Читателю доказательство приведённое выше может показаться незаконченным, поэтому распишем подробнее.\\
	Заметим, что для любой последовательности, какой бы она ни была, можно построить другую, для которой начальная будет последовательностью частичных сумм.\\
	Тогда из (\ref{13_th1}) следует, что предел последовательности равен пределу последовательности её средних арифметических.\\
	Положим $\ds a_n:=\ln S_n,\ b_n=\frac{a_1+\dots+a_n}n,\ \lim_{n\ra\i}a_n=\lim_{n\ra\i}b_n=\ln L$.
	\begin{multline*}
		\lim_{n\ra\i}\sqrt[n]{S_1\dots S_n}=\lim_{n\ra\i}\exp(\ln\sqrt[n]{S_1\dots S_n})=
		\lim_{n\ra\i}\exp\left(\frac1n\sum_{k=1}^n\ln S_k\right)=\\
		=\exp\left(\lim_{n\ra\i}\frac1n\sum_{k=1}^na_n\right)=\exp\left(\lim_{n\ra\i}b_n\right)=\exp(\ln L)=L
	\end{multline*}
\end{anote}
\begin{theorem}
	\[\lim_{n\ra\i}\frac{\ln p_1+\dots+\ln p_n}n=\ln L\ (-\i \text{ если }L=0)\]
\end{theorem}
\begin{proof}
	Было получено в замечании автора чуть выше.
\end{proof}
\begin{corollary}
	Если $\lim_{n\ra\i}\frac{a_{n+1}}{a_n}=l$, то $\lim_{n\ra\i}\sqrt[n]{a_n}=l$.
\end{corollary}
\begin{proof}
	Положим: $p_1:=a_1,\ p_2:=\frac{a_2}{a_1},\dots,\ p_n:=\frac{a_n}{a_{n-1}},\ \dots$.\\
	Тогда по следствию (\ref{13_cl1}):
	\[\lim_{n\ra\i}\sqrt[n]{p_1\dots p_n}=\lim_{n\ra\i}\sqrt[n]{a_n}=l\]
\end{proof}
\begin{theorem}\nf{(Критерий Коши сходимости числовых рядов)}\label{13_th8}\\
	Ряд $\ds\ss a_n$ сходится тогда и только тогда, когда:
	\[(\fa\eps>0)(\ex N\in\N)(\fa n>N)(\fa p\in\N)\ \l|\sum_{k=n+1}^{n+p}a_k\r|<\eps\]
\end{theorem}
\begin{proof}
	$\ss a_n$ cходится $\lra\ \ex\lim_{n\ra\i}S_n$, тогда по критерию Коши:
	\[(\fa\eps>0)(\ex N\in\N)(\fa n>N)(\fa p\in\N)\ \l|S_{n+p}-S_n\r|<\eps\]
\end{proof}
\begin{definition}
	Ряд $\ss a_n$ называется \textit{сходящимся абсолютно}, если $\sum_{n=1}^\i|a_n|$ сходится.
\end{definition}
\begin{theorem}
	Если ряд $\ss a_n$ сходится абсолютно, то он сходится.
\end{theorem}
\begin{proof}
	$\ss|a_n|$ сходится, тогда:
	\[(\fa\eps>0)(\ex N\in\N)(\fa n>N)(\fa p\in\N)\ \sum_{k=n+1}^{n+p}|a_k|<\eps\]
	\[\l|\sum_{k=n+1}^{n+p}a_k\r|\le\sum_{k=n+1}^{n+p}|a_k|<\eps\]
\end{proof}
\begin{definition}
	Если ряд сходится, но не абсолютно, то он сходится условно.
\end{definition}
\begin{theorem}
	Если $\ss a_n=\ss b+\ss c_n$, то:
	\begin{enumerate}
		\item Если $\ss b_n$ и $\ss c_n$ сходятся абсолютно, то $\ss a_n$ тоже сходится абсолютно.
		\item Если $\ss b_n$ условно сходится, а $\ss c_n$ сходится абсолютно, то $\ss a_n$ условно сходится, $\ss b_n$ и $\ss c_n$ можно поменять местами.
	\end{enumerate}
\end{theorem}
\begin{proof}
	\begin{enumerate}
		\item $|a_n|\le|b_n|+|c_n|$, тогда по признаку сравнения получим требуемое.
		\item От противного. Пусть $a_n$ сходится абсолютно, при этом $|b_n|\le|a_n|+|c_n|$, то есть, по признаку сравнения $b_n$ должно сходится абсолютно, противоречие с условием.
	\end{enumerate}
\end{proof}
\begin{theorem}\nf{(Признак Лейбница)}\label{13_th11}\\
	Если $b_n$ убывающая бесконечно малая последовательность, то знакоочередный ряд $\ss (-1)^{n+1}b_n$ сходится, причём для остатка $R_n:=\sum_{k=n+1}^\i(-1)^{k+1}b_k$ справедлива оценка $|R_n|\le b_{n+1}$
\end{theorem}
\begin{proof}
	Рассмотрим $S_{2n}=(b_1-b_2)+\dots+(b_{2n-1}-b_{2n})$
	\[S_{2(n+1)}=S_{2n+2}=S_{2n}+b_{2n+1}-b_{2n+2}\Ra\{S_{2n}\}_{n=1}^\i\text{ возрастает}\]
	\[S_{2n}=b_1-(b_2-b_3)-(b_4-b_5)-\dots-(b_{2n-2}-b_{2n-1})-b_{2n}\le b_1\]
	По т. Вейерштрасса: $\ex\lim_{n\ra\i}S_{2n}=S\in\R$
	\[S_{2n+1}=S_{2n}+b_{2n+1}\ra S\text{ при }n\ra\i\]
	\[R_n=(-1)^{n+1}(b_{n+1}-(b_{n+2}-b_{n+3})-\dots)\]
\end{proof}
\begin{theorem}\nf{(Признак Дирихле сходимости числовых рядов)}\label{13_th12}\\
	Если:
	\begin{enumerate}
		\item Частичные суммы $A_N:=\sum_{n=1}^Na_n$ ограничены.
		\item $\{b_n\}_{n=1}^\i$ монотонна и бесконечно мала.
	\end{enumerate}
	то $\ss a_nb_n$ сходится.
\end{theorem}
\begin{proof}
	По критерию Коши:
	\begin{multline*}
		\sum_{k=n+1}^{n+p}a_kb_k=\sum_{k=n+1}^{n+p}(A_k-A_{k-1})b_{n+1}+\dots+(A_{n+p}-A_{n+p-1})b_{n+p}=\\
		=A_{n+p}b_{n+p}-A_nb_{n+1}+\sum_{k=n+1}^{n+p-1}A_k(b_k-b_{k+1})
	\end{multline*}
	\[|A_N|\le M\ (\fa N\in\N)\]
	\[\l|\sum_{k=n+1}^{n+p}a_kb_k\r|\le M\l(|b_{n+1}|+|b_{n+1}|+\l|\sum_{k=n+1}^{n+p-1}(b_k-b_{k+1})\r|\r)\le2M(|b_{n+1}|+|b_{n+1}|)\]
	\[\lim_{n\ra\i}b_n=0\Ra(\fa\eps>0)(\ex N\in\N)(\fa n>N)\ |b_n|<\frac\eps{4M}\Ra(\fa p\in\N)\ |b_{n+p}|<\frac\eps{4M}\]
	Тогда по (\ref{13_th8}) получим требуемое.
\end{proof}
\begin{corollary}(\nf{Признак Абеля сходимости числовых рядов})\\
	Если:
	\begin{enumerate}
		\item Ряд $\ss a_n$ сходится.
		\item $\{b_n\}_{n=1}^\i$ монотонна и органичена.
	\end{enumerate}
	то $\ss a_nb_n$ cходится.
\end{corollary}
\begin{proof}
	Из $2$ пункта следует, что $\ex\lim_{n\ra\i}b_n=b\in\R$.\\
	\begin{equation*}
		\l.
		\begin{aligned}
			\t b_n:=b_n-b\Ra\{\t b_n\}\text{ монотонна и бесконечно мала}&\\
			\ss a_n\text{ сходится }\Ra\text{ частичные суммы ограничены}&
		\end{aligned}
		\r\}
		\Ra\ss a_n\t b_n\text{ сходится}
	\end{equation*}
	\[\ss a_nb_n=\ss a_n\t b_n+b\ss a_n\]
\end{proof}
\begin{examples}
	\[\ss\frac{\sin nx}{n^\al},\ \ss(1+\frac1n)^n\cdot\frac{\sin nx}{n^\al},\ x\ne\pi n,\ n\in\Z\]
	При $\al\le0$ расходится.\\
	От противного. Пусть $s=\lim_{n\ra\i}\sin nx$. Также пусть $c:=\lim_{n\ra\i}\cos nx$.
	\[\sin(n+1)x=\sin nx\cos x+\cos nx\sin x\]
	\[\cos(n+1)x=\cos nx\cos x-\sin x\sin x\]
	\begin{equation*}
		\begin{cases}
			s=s\cos x+c\sin x
			c=c\cos x-s\sin x
		\end{cases}
	\end{equation*}
	\[(s^2+c^2)=(s^2+c^2)\cos x=1\Ra\cos x=1\]
	Получили противоречие, значит $\nexists\lim_{n\ra\i}\sin nx$.
	\[\l|\sum_{n=1}^N\sin nx\cdot\frac{\sin\frac x2}{\sin\frac x2}\r|=\l|\sum_{n=1}^N\frac{\cos(n-\frac12)x-\cos(n+\frac12)x}{2\sin\frac x2}\r|=\l|\frac{\cos\frac x2-\cos(N+\frac12)x}{2\sin\frac x2}\r|\le\frac1{|\sin\frac x2|}\]
	$\{\frac1{n^\al}\}$ - монотонная и бесконечно мала при $\al>0$.\\
	$\ss\frac{\sin nx}{n^\al}$ сходится при $\al>0$.
	\[\l|\frac{\sin x}{n^\al}\r|\le\frac1{n^\al},\ \ss\frac1{n^\al}\text{ абсолютно сходится при }\al>1,\text{ условно при }\al\in(0,1]\]
	\begin{equation*}
		\l.
		\begin{aligned}
			\l|\frac{\sin nx}{n^\al}\r|\ge\frac{\sin^2nx}{n^\al}=\frac{1-\cos2\pi x}{2n^\al}&\\
			\ss\frac1{2n^\al}\text{ расходится при }\al\in(0,1]&\\
			\ss\frac{\cos 2\pi x}{n^\al}\text{ сходится при }\al>0
		\end{aligned}
		\r\}
		\Ra\ss\frac{\sin^2 n\al}{n^\al}\text{ расходится при }\al\in(0,1]
	\end{equation*}
\end{examples}
\begin{theorem}\nf{(О перестановках абсолютно сходящихся рядов)}\\
	Если ряд $\ss a_n$ сходится абсолютно, то $\fa$ биекции $n:\N\ra\N$ и последовательности $\{n_k\}_{k=1}^\i,\ n_k=n(k)$ ряд $\ss a_{n_k}$ сходится к той же сумме, что и $\ss a_n$.
\end{theorem}
\begin{proof}
	Положим $a_k^*:=a_{n_k}$\\
	\[\sum_{k=1}^N|a_k^*|\le\sum_{n=1}^\i|a_n|\Ra\sum_{k=1}^\i a_k^*\text{ абсолютно сходится по }(\ref{12_th2})\]
	\[S_n^*:=\sum_{k=1}^na_k^*,\ \sum_{k=1}^na_k=S_n\]
	\[S_m-S_n^*\ (\fa m\ge N)\text{, где }N\text{ - такой номер, что }\{a_1^*,\dots,a_n^*\}\subset\{a_1,\dots,a_N\}\]
	\[|S_m-S_n^*|\le\sum_{k=n+1}^\i|a_k^*|\underset{m\ra\i}{\longrightarrow}|S-S_n^*|\le\sum_{k=n+1}^\i|a_k^*|\underset{m\ra\i}{\longrightarrow}0\]
	Получили требуемое.
\end{proof}